\section{Next Steps}
  \label{sec:next}

  The programme has now established spectral admissibility on two finite
  subgroups of $\mathrm{SU}(2)$:
  \begin{center}
    \renewcommand{\arraystretch}{1.2}
    \begin{tabular}{lll}
      \hline
      Step & Group & Key result \\
      \hline
      1--3 & $Q_8$ & $\lambda_{\text{non-ab}} = 6 < \lambda_{\text{abel}} = 8$;
      admissibility ratio $\sqrt{4/3}$ \\
      4 & $Q_8$ & structural origin via Pauli tracelessness
      ($\mu_{\rho_4} = 0$) \\
      5 & $2I$ & four-level hierarchy; ratio $\sqrt{14/9}$;
      origin via $\varphi - 1/\varphi = 1$ \\
      \hline
    \end{tabular}
  \end{center}

  Both results are exact, non-perturbative, and derived from character
  theory alone.
  The hierarchy strengthens along the inclusion chain
  $Q_8 \subset 2I \subset \mathrm{SU}(2)$, and its algebraic origin is
  structural at each step.

  \paragraph{Step~6: Large Cayley graphs approaching $\mathrm{SU}(2)$.}

    The natural continuation is to test whether the spectral hierarchy
    persists and sharpens along a family of Cayley graphs whose size grows
    without bound while the geometry converges toward
    $\mathrm{SU}(2) \cong S^3$.
    The complete programme chain is
    \begin{equation}
      Q_8
      \;\longrightarrow\;
      2I
      \;\longrightarrow\;
      \mathrm{SL}(2,\mathbb{F}_p)
      \;\xrightarrow{\;p \to \infty\;}
      \mathrm{SU}(2),
      \label{eq:chain-full}
    \end{equation}
    where the final arrow denotes local spectral convergence of the Cayley
    graphs toward the $S^3$ Laplacian.

    Natural candidates are the Ramanujan graphs of Lubotzky, Phillips, and
    Sarnak~\cite{LPS1988} on $\mathrm{SL}(2,\mathbb{F}_p)$.
    These graphs have three properties that make them ideal for the
    programme:

    \begin{enumerate}
      \item \emph{Near-optimal spectral gap.}
      The non-trivial Laplacian eigenvalues satisfy
      $\lambda \ge d - 2\sqrt{d-1}$ (Alon--Boppana bound),
      where $d = p + 1 = |S|$.
      The admissibility envelope
      $A^{\max}_\rho = c_\chi/\sqrt{\lambda_\rho}$ therefore remains
      uniformly controlled as $p$ grows.

      \item \emph{Increasing size.}
      $|\mathrm{SL}(2,\mathbb{F}_p)| = p(p^2-1)$ provides a sequence of
      progressively finer discrete approximations to $S^3$.

      \item \emph{Full spin tower in the limit.}
      The irreducible representations of $\mathrm{SL}(2,\mathbb{F}_p)$
      include all half-integer spins
      $j = 0, \tfrac12, \ldots, \tfrac{p-1}{2}$,
      reproducing the complete spin spectrum of $\mathrm{SU}(2)$ as
      $p \to \infty$.
    \end{enumerate}

    Step~6 consists of four sub-tasks.

    \medskip
    \noindent\textbf{Step~6a (generating set).}
    The LPS construction provides an explicit symmetric set of $p+1$
    generators of $\mathrm{SL}(2,\mathbb{F}_p)$, defined by the solutions to
    $a^2+b^2+c^2+d^2 \equiv 1 \pmod p$ with $a$ odd and $b,c,d$ even,
    embedded via the standard quaternion-to-matrix map.
    This set plays the same structural role as the generators
    $\{\pm i,\pm j,\pm k\}$ of $Q_8$ and the order-10 elements of $2I$.

    \medskip
    \noindent\textbf{Step~6b (character computation).}
    For any Cayley graph generated by a union of conjugacy classes,
    the adjacency eigenvalue in the irreducible sector $\rho$ is
    \begin{equation}
      \mu_\rho
      =
      \frac{1}{\dim\rho}
      \sum_{C \subseteq S} |C|\,\chi_\rho(C).
      \label{eq:mu-SL2}
    \end{equation}
    The character tables of $\mathrm{SL}(2,\mathbb{F}_p)$ are classical,
    and the spin-$j$ sector corresponds to the representation
    $\mathrm{Sym}^{2j}(V_2)$, exactly as in the $Q_8$ and $2I$ cases.

    \medskip
    \noindent\textbf{Step~6c (admissibility envelope).}
    The Laplacian eigenvalue and admissibility window for the spin-$j$
    sector are
    \begin{equation}
      \lambda_j = (p+1) - \mu_j,
      \qquad
      A^{\max}_j = \frac{c_\chi}{\sqrt{\lambda_j}}.
      \label{eq:Amax-SL2}
    \end{equation}

    \medskip
    \noindent\textbf{Step~6d ($p\to\infty$ limit).}
    The key observable is the ratio
    \begin{equation}
      R_p
      =
      \frac{A^{\max}_{j=1/2}}{A^{\max}_{j=j_{\max}}}
      =
      \sqrt{\frac{\lambda_{j_{\max}}}{\lambda_{1/2}}},
      \qquad
      j_{\max} = \frac{p-1}{2}.
      \label{eq:ratio-p}
    \end{equation}

    The programme predicts $R_p>1$ for all finite $p$,
    meaning that the spin-$\tfrac12$ sector remains the most
    admissible sector along the sequence.

    If the ordering
    $\lambda_{1/2} < \lambda_1 < \cdots < \lambda_{j_{\max}}$
    becomes monotone for large $p$ and
    $\lim_{p\to\infty} R_p > 1$,
    this would establish
    \begin{equation}
      \text{bounded relational flux}
      \;\Longrightarrow\;
      \text{monotone low-spin survival}
      \quad
      \text{in the }\mathrm{SU}(2)\text{ limit.}
      \label{eq:conclusion}
    \end{equation}

    The results on $Q_8$ and $2I$ already provide the first two
    non-trivial instances of this mechanism.
    Their algebraic origins — Pauli tracelessness and the golden-ratio
    identity $\varphi - 1/\varphi = 1$ — are both instances of the same
    underlying principle: the generating set $S$ is positively correlated
    with the low-spin representations in character.
    This yields
    $\mu_{\text{low spin}} > \mu_{\text{high spin}}$
    and therefore
    $\lambda_{\text{low}} < \lambda_{\text{high}}$.

    Whether this structural correlation persists uniformly along the chain
    \eqref{eq:chain-full} is the central open question of Step~6.
